% Generated draft: requires mathematical and source review.
\documentclass[11pt]{article}
\usepackage[margin=25mm]{geometry}
\usepackage{fontspec}
\setmainfont{Latin Modern Roman}
\usepackage{amsmath,amssymb,mathtools,unicode-math}
\setmathfont{Latin Modern Math}
\usepackage{enumitem,needspace,xurl,hyperref}
\setlength{\parindent}{0pt}
\setlength{\parskip}{5pt}
\setlength{\emergencystretch}{3em}
\newcommand{\N}{\mathbb{N}}
\newcommand{\Z}{\mathbb{Z}}
\newcommand{\Q}{\mathbb{Q}}
\newcommand{\R}{\mathbb{R}}
\newcommand{\C}{\mathbb{C}}
\newcommand{\F}{\mathbb{F}}
\newcommand{\A}{\mathbb{A}}
\newcommand{\PP}{\mathbb{P}}
\newcommand{\E}{\mathbb{E}}
\newcommand{\eps}{\varepsilon}
\newcommand{\dd}{\,\mathrm{d}}
\begin{document}
\section*{UnsolvedMath ID: 30006229}
\textbf{Source problem code:} catalogue entry
\subsection{Definitions and mathematical statement}
Let \(M\) be a complete, connected, noncompact Riemannian manifold of dimension \(N\), with metric \(g\) and Riemannian volume measure \(\mu\). Write \(\mu(\Omega)\) for the quantity denoted by \(m(\Omega)\) in the source, distinguishing volume from the diffusion exponent \(m\). The supplied geometric hypotheses are
\[
\operatorname{Ric}_M\geq -(N-1)k\,g,
\qquad
\lambda_1(\Omega)\geq c\,\mu(\Omega)^{-2/N}
\]
for every nonempty relatively compact open set \(\Omega\subset M\). Under the usual Dirichlet interpretation,
\[
\lambda_1(\Omega)=\inf_{0\ne\varphi\in C_c^\infty(\Omega)}
\frac{\int_\Omega |\nabla\varphi|^2\,\dd\mu}
{\int_\Omega |\varphi|^2\,\dd\mu}.
\]
The record does not state the permitted ranges of \(k\) and \(c\), whether a boundary is excluded, or any boundary conditions; these require confirmation.

For \(s\in(0,1)\) and \(m>1\), consider the evolution equation
\[
\partial_t u=-(-\Delta_M)^s(u^m)
\qquad\text{on }(0,\infty)\times M.
\]
Here \(\Delta_M\) denotes the Laplace--Beltrami operator. The realization and definition of its fractional power are not supplied. The record also does not specify the admissible initial data, the sign convention for \(u^m\) when \(u<0\), or the definition of a weak dual solution.

\textbf{Requested assertions, subject to these unresolved specifications.}
\begin{enumerate}
\item Establish uniqueness of weak dual solutions: two solutions in the intended admissible class with the same initial datum must agree. The class, initial trace, and sense of agreement must be specified before this becomes a fully quantified assertion.
\item For each pole \(y\in M\), construct a fundamental solution with initial datum \(\delta_y\), the unit Dirac measure at \(y\). The distributional meaning of this initial datum is
\[
\lim_{t\downarrow0}\int_M U_y(t,x)\varphi(x)\,\dd\mu(x)=\varphi(y)
\qquad\text{for every }\varphi\in C_c^\infty(M).
\]
The solution class, additional trace or integrability requirements, and any mass conservation requirement are unspecified.
\item Establish pointwise bounds for general solutions in the intended class. No definite bound, dependence of constants, spatial regime, or meaning of pointwise evaluation is provided, so this target cannot be made into a specific inequality from the supplied record alone.
\end{enumerate}
The original excerpt also mentions existence of a weak dual solution for data in a space apparently written \(L^1_{G^s}(M)\), and smoothing estimates controlling \(\|u(t)\|_{L^\infty(M)}\) in terms of a norm of the initial datum. That space is undefined in the record. Its example estimate is corrupted: it contains \(\vartheta_1\), powers involving \(2s\vartheta_1\) and \(N\vartheta_1\), a maximum symbol, and the quantifier \(t>0\), but their complete arrangement and parameter definitions cannot be recovered reliably. The excerpt does not clearly distinguish whether these additional bullets describe established results or further requested assertions.
\subsection{Short English statement}
For fractional porous medium diffusion on a complete noncompact manifold satisfying the supplied curvature and eigenvalue inequalities, the requested goals are uniqueness of weak dual solutions, existence of solutions starting from a Dirac mass, and pointwise estimates for general solutions. Essential solution definitions, data classes, operator conventions, and the desired estimates are missing. The supplied open status and verification labels are curation claims; this statement-only draft requires human mathematical review.
\subsection{Sources}
\begin{itemize}
\item Supplied local UnsolvedMath registry record, permanent ID \(30006229\), title: \emph{Uniqueness and Fundamental Solutions for Fractional Diffusion}; source problem code \texttt{catalogue entry}. The record attributes the problem to the Oberwolfach Reports Open Problems collection and gives the proposed year as 2025; no individual proposer is supplied.
\item \emph{Mini-Workshop: Hardy Inequalities in Discrete and Continuum Settings} (2025), Oberwolfach Reports, report 14/2025, pp.\ 579--614, DOI \texttt{10.4171/owr/2025/14}. Supplied source and literature link: \url{https://doi.org/10.4171/owr/2025/14}. No problem-specific page or named contribution is supplied.
\item Supplied literature assessment, checked \texttt{2026-08-22}: the report is said to leave uniqueness, fundamental solutions with Dirac initial data, and pointwise bounds for general solutions open; the supplied assessment says that no later source was located resolving these claims in the stated generality. This is supplied dated curation, not an independently verified assessment. The record labels the problem \texttt{open} and the statement \texttt{corrected\_verified}, and says that the reconstructed statement was checked against the official report. Those labels and verification claims are preserved as supplied provenance, not adopted as independent verification.
\end{itemize}
\end{document}
